\section{Supporting comparisons under a fixed calibration}
\label{sec:conference-comparisons}

In this section, we fix $\eta_U = 4h$ and $\eta_S = h$. Both rules then have
diagonal drift $h$ in the common coordinate. Under independence, this gives the
common coefficient the same mean trajectory for both rules, although dependence
generally breaks the match. We specify $h$ in advance and use it for both
methods.

\subsection{Learning the rare signal versus fitting noise}

Under independence, SD scaling changes the contraction factor for the mean rare
coefficient error from $1 - 4hb$ to $1 - h$. At the same time, it increases the
observation-noise variance added to that coefficient at each update from
$16h^2b\sigma^2/B$ to $h^2\sigma^2/(bB)$. Standardization therefore speeds up
learning while introducing more noise. Random feature sampling contributes
additional variance to both procedures.

The diagonal error moments satisfy a nonnegative two-dimensional recurrence. By
omitting the terms that couple the two coordinates, we obtain a scalar
approximation to rare risk with an explicit nonnegative remainder. Let $F$ and
$N$ denote the noiseless and unit-noise parts of the SD-minus-U risk difference.
The approximation gives bounds
$F_- \leq F \leq F_+$ and $N_- \leq N \leq N_+$. When $F_+ < 0 < N_-$, the
noise-variance threshold $\tau = -F/N$ satisfies
\begin{equation}
 \frac{-F_+}{N_+} \leq \tau
                  \leq \frac{-F_-}{N_-}. \label{eq:conference-threshold}
\end{equation}
The remainder accounts for the dependence that the sampled updates introduce
between coefficient errors, even when the features are independent. We derive
the approximation and define its terms in Appendix~\ref{sec:comparison-proofs}.

\begin{theorem}[Stopping and batch reversal]
 \label{thm:conference-c1}
 In the binary model with $\ell = 0$, suppose
 \begin{gather*}
  q \in[0.019,0.021] ,\quad h \in[0.0019,0.0021] ,               \\
  \beta_C^*,\beta_R^* \in[0.95,1.05] ,\quad\sigma^2 \in[1.5,5] .
 \end{gather*}
 Both U and S are mean-square stable for $B \in \{1,10\}$. Their rare-risk
 contrast $\Delta_{B,k} = R_{R,k}^S - R_{R,k}^U$ lies in the intervals in
 Table~\ref{tab:conference-c1} throughout this parameter box.
\end{theorem}

\begin{table}[t]
 \centering
 \begin{tabular}{rrrr}
  \toprule
  $B$ &    $k$ & $M = kB$ &    $\Delta_{B,k}$ \\
  \midrule
    1 &  1,000 &    1,000 &   $[-.353,-.160]$ \\
    1 & 10,000 &   10,000 &    $[.0021,.125]$ \\
   10 &  1,000 &   10,000 &   $[-.386,-.286]$ \\
   10 & 10,000 &  100,000 & $[-.0281,-.0016]$ \\
  \bottomrule
 \end{tabular}
 \caption{
  Bounds on the rare-risk difference in Theorem~\ref{thm:conference-c1}, valid
  throughout the parameter box. Negative values favor SD scaling. The intervals
  enclose the exact risk differences and do not represent Monte Carlo
  uncertainty.
 }
 \label{tab:conference-c1}
\end{table}

At $B = 1$, the ordering reverses as training continues. With the same budget
$M = 10{,}000$, it also differs between $B = 1$ and $B = 10$. These are
comparisons between normalization rules at each batch size and budget. They do
not establish which batch size is better for either rule.

For the anchor $q = .02$, $h = .002$, $\beta^* = (1,1)^\top$, and
$\sigma^2 = 3$, the early $B = 1$ rare risk is $.356678$ for U and $.090407$ for
S. Squared bias contributes $.353207$ and $.011040$, respectively, so the early
advantage comes mainly from learning the rare signal more quickly. By
$k = 10{,}000$, U has learned enough of that signal to reach risk $.028145$,
whereas S has risk $.078788$, almost entirely due to observation noise. The
exact calculation includes variation from feature sampling at every step.

\subsection{Aggregate and rare-presence objectives}

Dependence can make the two prediction objectives favor different rules. Let
$\Delta V = V_k^S - V_k^U$. Direct substitution into \eqref{eq:conference-risks}
gives
\begin{equation}
 \Delta_R - \Delta_A = \frac{1 - 2q}{2}\bigl\{(1 - q)\Delta V_{RR} + \ell\Delta V_{CR}\bigr\}, \label{eq:conference-objectives}
\end{equation}
where $\Delta_A$ and $\Delta_R$ denote the aggregate and rare-risk differences.
The cross moment is the sum of the product of the coefficient biases and the
covariance of their errors.

\begin{theorem}[Objective disagreement]
 \label{thm:conference-c2}
 Take $q = .02$, $\ell \in[.9,1]$, $\beta^* = (2,1)^\top$, $h = .01$, $B = 1$,
 $k = M = 100$, and $\sigma^2 \in[3.8,4.1]$. Both U and S are mean-square
 stable, and uniformly over these intervals,
 \begin{align*}
  .00077 & \leq R_k^S - R_k^U \leq .00194,         \\
  -.0983 & \leq R_{R,k}^S - R_{R,k}^U \leq -.0233.
 \end{align*}
 Moreover, $V_{CC,k}^S > V_{CC,k}^U$ and $V_{RR,k}^S > V_{RR,k}^U$.
\end{theorem}

SD scaling increases each coefficient's mean squared error but improves
rare-presence prediction. At $\ell = 1$ and $\sigma^2 = 4$,
\[
 \Delta V \approx \begin{pmatrix}
                   .022743  & -.615146 \\
                   -.615146 & .470335
                  \end{pmatrix}.
\]
The reduction in the cross moment offsets the larger individual errors when
predicting for the rare subgroup. Aggregate risk gives this moment weight
$.01$, whereas rare risk gives it weight $.49$. Looking only at the individual
coefficient errors would therefore miss the improvement in rare prediction.

We prove both theorems using rational interval arithmetic with outward rounding
over the full parameter boxes. Appendix~\ref{sec:comparison-proofs} gives the
recurrences and enclosure arguments. The results apply to the specified
training procedures, whereas Theorem~\ref{thm:conference-predictive} compares
risks after optimizing over all constant rates and deterministic stops. Adaptive
stopping and the choice of an optimal normalization remain outside these
results. The identity relating the two risk objectives is elementary.
